\documentclass{article}
\usepackage[utf8]{inputenc}
\usepackage[T1]{fontenc}
\usepackage{geometry}
\geometry{
    a4paper,
    left=2.5cm,
    right=2.5cm,
    top=2.5cm,
    bottom=2.5cm
}
\usepackage{graphicx}
\usepackage{float}
\usepackage{hyperref}
\usepackage[english]{babel}
\usepackage{amsmath}
\usepackage{amssymb}
\usepackage{array}
\usepackage{booktabs}
\usepackage{caption}
\usepackage{subcaption}
\usepackage{enumitem}
\usepackage{fancyhdr}
\usepackage{titlesec}
\usepackage{url}
\usepackage{textcomp}
\usepackage{xcolor}
\usepackage{listings}
\usepackage{multirow}
\usepackage{longtable}
\usepackage{tabularx}
\usepackage{times}

% Set up page style
\pagestyle{fancy}
\fancyhf{}
\rfoot{\thepage}
\renewcommand{\headrulewidth}{0pt}

% Title formatting
\titleformat{\chapter}[display]{\normalfont\bfseries}{\chaptertitle\thechapter}{0pt}{\Huge}
\titlespacing*{\chapter}{0pt}{50pt}{40pt}

% Document info
\title{Template Document}
\author{Academic Support Agent}
\date{\today}

\begin{document}
\maketitle

\begin{titlepage}

\noindent
\Huge Reporte de laboratorio:\\ nombre del experimento
\epigraph{Resumen:\\ \lipsum[1] }%

\null\vfill
\vspace*{1cm}
\noindent
\hfill
\begin{minipage}{0.35\linewidth}
    \begin{flushright}
        \printauthor
    \end{flushright}
\end{minipage}
%
\begin{minipage}{0.02\linewidth}
    \rule{1pt}{125pt}
\end{minipage}
\titlepagedecoration
\end{titlepage}

\newpage

\section{Introducci n / Marco te rico}

\subsection{Subtema}
El subtema trata sobre lo que significa ser una ecuaci n numerada como la siguiente:
\begin{equation}\label{sch}
       \frac{\hbar^2}{2m}\nabla^2\Psi + V(\mathbf{r})\Psi = -i\hbar \frac{\partial\Psi}{\partial t}
\end{equation}
pero tambi n podemos tener ecuaciones no numeradas como
\begin{equation*}
    ds^2 = c^2 dt^2 \left( \frac{d\sigma^2}{1-k\sigma^2} 
    + \sigma^2\left[ d\theta^2 
    + \sin^2\theta d\phi^2 \right] \right)
\end{equation*}
y podemos hacer referencia a cualquiera de ellas usando su etiqueta (label). Por ejemplo, 
refiri ndose a la ec. (\ref{sch}).

Tambi n podemos tener ecuaciones agrupadas y alineadas
\begin{equation}
\begin{aligned} 
2x - 5y &=  8 \\ 
3x + 9y &=  -12.
\end{aligned}
\end{equation}

o ecuaciones agrupadas por un corchete
\begin{equation}
\begin{cases}
\begin{aligned} 
2x - 5y &=  8 \\ 
3x + 9y &=  -12.
\end{aligned}
\end{cases}
\end{equation}

\subsubsection{Subsubtema }

\section{Procedimiento experimental}

\subsection{Material}

\begin{enumerate}
    \item Objeto 1
    \begin{itemize}
        \item Parte a
        \item Parte b
    \end{itemize}
    \item Objeto 2
    \item Objeto 3
\end{enumerate}

\subsection{Experimento 1}

\subsection{Experimento 2}

\section{Resultados}

\subsection{Resultados experimento 1}

\subsection{Resultados experimento 2}

\section{An lisis y discusi n}

\subsection{Cuestionario}

\section{Conclusiones}

\section{Bibliografia}
\end{document}